% =====================================================================
% Operator.tex
% Math operators -- behave like \sin / \log etc.
% Part of the definitions-latex package.
%
% Required packages: amsmath, amssymb
% =====================================================================

% ---------------------------------------------------------------------
% Generic
% ---------------------------------------------------------------------
\DeclareMathOperator{\cte}{cte}            % constant
\DeclareMathOperator{\atanh}{atanh}        % inverse hyperbolic tangent
\DeclareMathOperator{\sgn}{sgn}            % sign function
\DeclareMathOperator{\extr}{extr}          % extremum
\DeclareMathOperator{\rank}{rank}          % matrix rank
\DeclareMathOperator{\multalg}{multalg}    % algebraic multiplicity
\DeclareMathOperator{\ima}{i}              % imaginary unit (operator-shaped)

% ---------------------------------------------------------------------
% Probability and statistics
% ---------------------------------------------------------------------
\DeclareMathOperator{\esp}{\mathbb{E}}     % expectation
\DeclareMathOperator{\var}{Var}            % variance
\DeclareMathOperator{\std}{\sigma}         % standard deviation
\DeclareMathOperator{\med}{median}         % median
\DeclareMathOperator{\iqr}{IQR}            % interquartile range
\DeclareMathOperator{\prob}{\mathbb{P}}    % probability measure
\DeclareMathOperator{\R}{\mathbb{R}}       % real numbers

% ---------------------------------------------------------------------
% Error / regression metrics
% ---------------------------------------------------------------------
\DeclareMathOperator{\rmse}{RMSE}
\DeclareMathOperator{\nrmse}{NRMSE}
\DeclareMathOperator{\rms}{RMS}
\DeclareMathOperator{\mse}{MSE}
\DeclareMathOperator{\mstr}{MSTR}

% ---------------------------------------------------------------------
% Diagnostic / inferential
% ---------------------------------------------------------------------
\DeclareMathOperator{\vif}{VIF}
\DeclareMathOperator{\pvalue}{\mathit{p}\text{-value}}    % FIXED: text hyphen, not minus

% ---------------------------------------------------------------------
% Mechanics
% ---------------------------------------------------------------------
\DeclareMathOperator{\ham}{\mathcal{H}}    % Hamiltonian